For the full-Haar pure $SU(2)$ Hamiltonian, the electric Casimir is unbounded
on each link, while a plaquette interaction has norm $\beta=2/(g^2a)$.
Four plaquettes meet a link and at most sixteen overlap one plaquette.
The interaction-picture commutator bound therefore gives, for disjoint
unit-norm observables at link-graph distance $d$,
\[
 \|[\tau_t(A),B]\|\leq
 \min\left\{2,\;2|X|\sum_{n\geq d}
                  \frac{(32\beta|t|)^n}{n!}\right\}.
\]
A polynomial boundary count and the same factorial tail prove norm
convergence of local finite-volume dynamics at fixed spacing, with every
Haar representation retained. At $g^2=32$, $a=t=1$, a four-link observable
with buffer $32$ has a certified boundary error below $3.14\times10^{-20}$.
This does not assign a numerical gap at that coupling. At fixed $g$, the
Hamiltonian and its nonzero gaps scale as $1/a$; a finite physical mass needs
a controlled changing-coupling or energy-unit trajectory. Scalar reference
subtraction cannot change a gap. Infinite volume and continuous Hamiltonian
time supply neither a four-dimensional spatial-continuum construction nor
a Clay mass-gap result.

The propagation bound is independent of the compact gauge group. For a
compact connected semisimple group with first Casimir eigenvalue $c_G>0$,
the same electric-star proof has local perturbation norm at most
$48/[c_G(g^2)^2]$ and retains an existential fixed-spacing uniform gap.
For $SU(N)$, positivity of the inverse Cartan matrix gives
$c_G=(N^2-1)/(2N)$; nonzero fundamental Wilson variance supplies physical
excitations in a sufficiently small perturbation interval. This extends the
fixed-spacing argument across the gauge groups without supplying a
spatial-continuum construction or a numerical coupling threshold.
